\section{Local counts, growth tails, and an exact fallback}\label{sec:count}

This section develops the probabilistic and algorithmic tools shared by the
quadratic, sparse, and recourse results. Every statement concerns a fixed
base function and an independent random linear perturbation of it. No
statement uses convexity, differentiability, uniqueness of an optimizer, or
quadratic growth of the base problem unless that property is listed among
its hypotheses.

Throughout, $I$ denotes the bit length of the rational base data before any
perturbation is sampled. It includes the noise width $\sigma$ and every
supplied decomposition or certificate. A perturbation is drawn once, from a
law that the algorithm computes from the base data before sampling. No
algorithm in this paper discards, resamples, or conditions on a draw, and
every algorithm is correct on every draw. Probability enters only the bounds
on work.

The section has four parts. Theorem~\ref{thm:count:local} bounds the
expected number of grid nodes that survive comparisons with their grid
neighbors. Together with curvature-corrected refinement
(Theorem~\ref{thm:count:cells}), it bounds the expected number of cells that
a search retains at every refinement level, uniformly in the level.
Section~\ref{sec:count:laws} transfers such bounds from continuous proxy
laws to finite rational laws that a Turing machine can sample.
Theorem~\ref{thm:count:growth-tail} and Lemma~\ref{lem:count:finite-tails}
bound the probability of small global quadratic growth and of small active
gradients. Theorem~\ref{thm:count:fallback} gives an exact algorithm for
every draw of a fixed-degree polynomial problem. Its cost contains an
exponential factor that depends only on the base data, and
Lemma~\ref{lem:count:rare-fallback} shows how a rare exceptional event pays
for that factor. Section~\ref{sec:count:template} summarizes how the later
sections combine these parts. Proofs are in Appendix~\ref{app:count}.

\subsection{Upper curvature and local comparison events}\label{sec:count:local}

Let $X=\prod_{i=1}^k[a_i,b_i]$ be a box with $a_i<b_i$ and widths
$w_i=b_i-a_i$. The mesh statements use $k\ge1$. Fixed coordinates are
held at their values, and a singleton box is evaluated directly before a
width maximum, mesh, or accuracy threshold is formed.
For a function $f:X\to\R$ and a coefficient vector
$c\in\R^k$ write
\[
 F_c(x)=f(x)+c^{\mathsf T} x .
\]

\begin{definition}[Upper coordinate curvature]\label{def:count:curv}
A function $f:X\to\R$ has \emph{upper coordinate curvature at most}
$L=(L_1,\dots,L_k)\in[0,\infty)^k$ if, for every $i$ and every $x\in X$, the
function $t\mapsto f(x+te_i)-L_it^2/2$ is concave on the interval
$\{t:x+te_i\in X\}$.
\end{definition}

The property holds when $f-\frac12\sum_iL_ix_i^2$ is concave. It is
preserved by pointwise infima: if every $f_\omega$ has upper coordinate
curvature at most $L$ and $\inf_\omega f_\omega$ is finite, then
$\inf_\omega f_\omega$ has upper coordinate curvature at most $L$, because an
infimum of concave functions is concave. The auxiliary value functions of
Section~\ref{sec:qp} are of this kind. Downward kinks, discontinuous
derivatives, and arbitrarily large negative curvature are allowed. Only an
upper bound on curvature is used.

\begin{lemma}[Neighbor comparisons confine one coefficient]\label{lem:count:interval}
Let $f:X\to\R$ be arbitrary, $v\in X$, $h>0$, $\delta\ge0$, and suppose
that $v\pm he_i\in X$. If
\[
 F_c(v)\le F_c(v+he_i)+\delta
 \qquad\text{and}\qquad
 F_c(v)\le F_c(v-he_i)+\delta ,
\]
then $c_i$ belongs to the interval
\[
 I_i(v,h,\delta)=\Bigl[\frac{f(v)-f(v+he_i)-\delta}{h},\
 \frac{f(v-he_i)-f(v)+\delta}{h}\Bigr].
\]
This interval depends on $f$, $v$, $h$, and $\delta$, but not on $c$. If $f$
has upper coordinate curvature at most $L$, its length is at most
$L_ih+2\delta/h$; an interval with negative length is empty.
\end{lemma}

\begin{definition}[Grids, local events, and interval concentration]\label{def:count:local-event}
For integers $m_i\ge1$ put $h_i=w_i/m_i$,
$G_i=\{a_i+\ell h_i:0\le\ell\le m_i\}$, and $G=\prod_iG_i$. The
\emph{interior coordinates} of a node $v\in G$ are
$Q_v=\{i:a_i<v_i<b_i\}$. For $\delta\ge0$, the \emph{$\delta$-local event}
$E_v^\delta$ is the set of coefficient vectors $c$ for which
$F_c(v)\le F_c(v\pm h_ie_i)+\delta$ for all $i\in Q_v$ and both signs. The
\emph{interval concentration} of a real random variable $Y$ is
\[
 \pi_Y(\ell)=\sup\{\Prob(Y\in J): J\text{ an interval of length at most }\ell\}.
\]
\end{definition}

A node with $F_c(v)\le\inf_XF_c+\delta$, or with
$F_c(v)\le\min_GF_c+\delta$, satisfies $E_v^\delta$, because its grid
neighbors are feasible competitors. Three concentration bounds recur. If $Y$
has a density bounded by $\phi$, then $\pi_Y(\ell)\le\phi\ell$. If $Y$ has the
grid law $U_{\sigma,M}$ of Definition~\ref{def:model:grid}, uniform on
$\{-\sigma+2\sigma j/(M-1):0\le j<M\}$, then
$\pi_Y(\ell)\le\ell/(2\sigma)+1/M$ by \eqref{eq:model:interval}.
If the distribution function of $Y$ differs everywhere by at most
$\delta_{\rm K}$ from that of a law with density bounded by $\phi$, then
$\pi_Y(\ell)\le\phi\ell+2\delta_{\rm K}$.

\begin{theorem}[Expected local counts]\label{thm:count:local}
Let $f:X\to\R$ have upper coordinate curvature at most $L$, let $G$ be the
grid of Definition~\ref{def:count:local-event}, and let the coordinates of
$c$ be independent, with interval concentrations $\pi_1,\dots,\pi_k$. Then
for every $\delta\ge0$
\[
 \E\,\#\{v\in G: c\in E_v^\delta\}
 \le\prod_{i=1}^k\Bigl[2+(m_i-1)\min\bigl\{1,\pi_i(L_ih_i+2\delta/h_i)\bigr\}\Bigr].
\]
The same bound holds conditionally on any random element that is
independent of $c$, if $f$ is allowed to depend on that element.
\end{theorem}

The bound concerns a deterministic grid. An adaptive search visits a subset
of its nodes that depends on the draw, but every node that the search keeps
for a probabilistic reason will be shown to satisfy a local event. The
expectation is therefore taken over a fixed family of events, and no
independence between the draw and the adaptively selected cells is used.
The comparisons with grid neighbors make the other coefficients cancel, which
is why independence across coordinates multiplies the per-coordinate bounds.

\begin{definition}[Balanced nested meshes]\label{def:count:mesh}
Let $\rho_1,\dots,\rho_k>0$ be \emph{mesh scales}. Put
$S_\rho=\max_iw_i/\rho_i$ and $h_j=S_\rho2^{-j}$ for $j\ge0$. At level $j$,
coordinate $i$ is divided into $m_{ij}$ equal parts, where $m_{ij}$ is the
least power of two with $w_i/m_{ij}\le\rho_ih_j$, and $h_{ij}=w_i/m_{ij}$.
The level-$j$ grid $G_j$ and the level-$j$ cells are the corresponding
tensor grid and the boxes between consecutive grid nodes. For a curvature
vector $L\in(0,\infty)^k$ put
\[
 E_j=\frac18\sum_{i=1}^kL_ih_{ij}^2,
 \qquad
 c_{\rm bal}=\max_{i,l}\frac{L_l\rho_l^2}{L_i\rho_i^2},
 \qquad
 C_k=1+2c_{\rm bal}k .
\]
The \emph{isotropic} choice $\rho_i=1$, for a common curvature $L_i=L$, gives
$c_{\rm bal}=1$ and $C_k=1+2k$.
\end{definition}

\begin{corollary}[Level-uniform counts]\label{cor:count:levels}
In Definition~\ref{def:count:mesh}, the partitions are nested, level $0$
consists of the single cell $X$, and $m_{ij}\le2^j$. Every coordinate with
$m_{ij}\ge2$ satisfies $\rho_ih_j/2<h_{ij}\le\rho_ih_j$ and
\[
 L_ih_{ij}+4E_j/h_{ij}\le C_kL_ih_{ij}.
\]
Consequently, if $f$ has upper coordinate curvature at most $L$ and the
coordinates of $c$ are independent, then at every level $j$
\[
 \E\,\#\{v\in G_j:c\in E_v^{2E_j}\}\le
 \begin{cases}
 \prod_{i=1}^k\bigl[2+C_kL_i\phi_iw_i\bigr]
 &\text{if $c_i$ has a density bounded by $\phi_i$,}\\[3pt]
 \prod_{i=1}^k\bigl[3+C_kL_iw_i/(2\sigma_i)\bigr]
 &\text{if $c_i\sim U_{\sigma_i,M}$ and $M\ge m_{ij}$ for all $i$.}
 \end{cases}
\]
\end{corollary}

A coordinate that is not yet refined has no interior node. It contributes
the factor $2$, with no dependence on its width, however small. The
construction divides every coordinate interval into equal parts. Clipping a
uniform step at the end of an interval would create a short final gap $g$,
and Lemma~\ref{lem:count:interval} would then produce a term $\delta/g$ that
the count cannot control. The balance constant $c_{\rm bal}$ matters only
for anisotropic curvature; it is used with $c_{\rm bal}<4$ in
Section~\ref{sec:qp:separable}.

\subsection{Curvature-corrected refinement and the role of closure}\label{sec:count:search}

\begin{lemma}[Corrected corners]\label{lem:count:rounding}
Let $f$ have upper coordinate curvature at most $L$, and let
$C=\prod_i[l_i,l_i+\eta_i]\subseteq X$ be a cell. For every $c$ and every
$x\in C$,
\[
 \min_{v\text{ corner of }C}F_c(v)\le F_c(x)+\frac18\sum_{i=1}^kL_i\eta_i^2 .
\]
\end{lemma}

\begin{definition}[Corrected-corner search]\label{def:count:search}
Fix balanced nested meshes on $X$, a level cap $J$, and exact evaluation of
$F_c$. Level $0$ processes the cell $X$. Level $j\ge1$ processes the
children of the cells retained at level $j-1$. For every processed cell $C$,
evaluate $F_c$ at all corners of $C$ and put
$\underline F(C)=\min_{\rm corners}F_c-E_j$. Optionally, a \emph{closure
step} computes $\min_CF_c$ exactly together with a point attaining it; such
a cell is \emph{closed}. After all cells of level $j$ have been processed,
let $U_j$ be the least value of $F_c$ found so far, at corners or at closure
points, stored together with a point attaining it. Retain exactly the
processed cells that are not closed and satisfy $\underline F(C)\le U_j$.
Stop when no cell is retained or when level $J$ has been processed.
\end{definition}

\begin{theorem}[Corrected-corner search]\label{thm:count:cells}
Let $f:X\to\R$ be lower semicontinuous with upper coordinate curvature at
most $L$, and put $F^*=\min_XF_c$. For every coefficient vector $c$ the
search of Definition~\ref{def:count:search} has the following properties.
\begin{enumerate}[label=(\roman*)]
\item $\underline F(C)\le\min_CF_c$ for every processed cell $C$.
\item If no cell is retained after level $j$, then $U_j=F^*$ and the stored
 point is a global minimizer.
\item $U_j\le F^*+E_j$.
\item Every cell retained at level $j$ has a corner $v\in G_j$ with
 $F_c(v)\le F^*+2E_j$. Hence $c\in E_v^{2E_j}$, and the number of retained
 cells is at most $2^k\#\{v\in G_j:c\in E_v^{2E_j}\}$.
\item The number $\underline U_j=\min\{U_j,\min_{C\text{ retained}}\underline F(C)\}$
 satisfies $\underline U_j\le F^*\le U_j$ and $U_j-\underline U_j\le E_j$,
 with the minimum over an empty retained set defined as $+\infty$.
\end{enumerate}
Consequently, under the hypotheses of Corollary~\ref{cor:count:levels}, the
expected number of cells retained at level $j$ is at most $2^k$ times the
corresponding bound of that corollary. This holds with or without closure
steps.
\end{theorem}

Item (iv) is the bridge between the adaptive search and the deterministic
count: a cell is retained for a probabilistic reason only through one of its
corners, which is a nearly optimal node of a fixed grid.

\begin{corollary}[Approximation with logarithmic accuracy dependence]\label{cor:count:approx}
Use isotropic meshes with common curvature $L>0$, let $s=\max_iw_i$, fix
$\varepsilon>0$, and put
$J=\max\{0,\lceil\frac12\log_2(kLs^2/(8\varepsilon))\rceil\}$. Without
closure steps, the search, stopped after level $J$ or earlier if no cell is
retained, returns a point $x$ and a number $\underline U$ with
$\underline U\le F^*\le F_c(x)$ and $F_c(x)-\underline U\le\varepsilon$, for
every $c$. If the coordinates of $c$ are independent with densities bounded
by $\phi_i$, the expected number of evaluations of $F_c$ is at most
$2^k+8^k\mathcal HJ$, where $\mathcal H=\prod_i[2+(1+2k)L\phi_iw_i]$. If
instead $c_i\sim U_{\sigma,M}$ with $M\ge2^J$, the same holds with
$\mathcal H=\prod_i[3+(1+2k)Lw_i/(2\sigma)]$.
\end{corollary}

For fixed $k$ and numerical bounds on $L\phi_iw_i$, the expected number of
evaluations is logarithmic in $1/\varepsilon$. The bound counts exact
evaluations of the sampled objective. Their bit cost must be added, and the
conclusion concerns the sampled objective, not the unperturbed one. The
finite law in the corollary is chosen for the requested accuracy. A fixed
finite law does not support the same bound at arbitrarily fine levels.

\begin{example}[A fixed atom defeats level-uniform counts]\label{ex:count:atom}
Let $k=1$ and $X=[0,1]$, and let the law of $c$ have an atom $c_0$ with
$\Prob(c=c_0)=p_0>0$. For $f(x)=-c_0x$, which has upper coordinate curvature at
most $L$ for every $L>0$, the draw $c=c_0$ makes $F_c\equiv0$. On this draw every node
of every level is a global minimizer, every cell has
$\underline F(C)=-E_j<0=U_j$, and all $2^j$ cells of level $j$ are retained.
The expected number of retained cells is therefore at least $p_02^j$, which
grows without bound.
\end{example}

The example limits this particular search: under a fixed finite law, the
cells retained by indefinitely continued refinement cannot be counted
uniformly in the level. It does not show that every exact algorithm needs a
rare fallback. A finite law supports Corollary~\ref{cor:count:levels} only up
to a level $J$ with $M\ge m_{iJ}$ (or, for transferred bounds, up to a node
budget fixed in advance), so an exact algorithm built on this search stops
refining at a level computed from the base data and finishes the sampled
problem in another way. Most results below use \emph{closure steps}, which
certify the exact minimum of a whole cell from a structural formula,
together with an exact \emph{fallback} on the same draw for the draws that
are still unresolved at level $J$. The lattice theorem
(Theorem~\ref{thm:int:lowrank}) instead stops at a predetermined level at
which the certified gap is below the spacing of the attainable objective
values, and the strong-noise theorem (Theorem~\ref{thm:int:strong-field})
does not use this search; neither needs a rare fallback. In all cases atoms,
ties, and flat optimal faces are handled on every draw; they are never
excluded by a genericity assumption.

\subsection{Finite rational laws and continuous proxies}\label{sec:count:laws}

A Turing machine with fair random bits can sample laws with infinitely many
rational atoms if its running time is unbounded; for example, counting the
failures before the first success gives a geometric law on the nonnegative
integers. The samplers in this paper are not of this kind: each uses at most
a prescribed number of fair bits and bit operations in the worst case, and a
sampler with this property has a finite law with rational atoms. The analysis
often computes a probability under a continuous proxy law and then transfers
it to the finite law that the algorithm actually samples. Two families of
finite laws are used.

\begin{definition}[Finite laws and Kolmogorov distance]\label{def:count:laws}
The grid law $U_{\sigma,M}$ is that of Definition~\ref{def:model:grid}. For an
integer $b\ge1$ and rational $\sigma>0$,
$\mathcal G_{b,\sigma}$ is the law of $\sigma Z_b$, where $Z_b$ is the
output of the sampler of Lemma~\ref{lem:count:gauss-sampler}. The
\emph{Kolmogorov distance} of laws $\mu,\nu$ on $\R$ is
$d_{\rm K}(\mu,\nu)=\sup_{t\in\R}|\mu(-\infty,t]-\nu(-\infty,t]|$.
\end{definition}

\begin{lemma}[Product replacement]\label{lem:count:transfer}
\begin{enumerate}[label=(\alph*)]
\item The uniform grid law satisfies
 $d_{\rm K}(U_{\sigma,M},\mathrm{Unif}[-\sigma,\sigma])\le1/M$, and
 the Gaussian-like law satisfies
 $d_{\rm K}(\mathcal G_{b,\sigma},N(0,\sigma^2))\le2^{-b}$.
\item Let $\mu_1,\dots,\mu_n$ and $\nu_1,\dots,\nu_n$ be laws on $\R$ with
 $d_{\rm K}(\mu_i,\nu_i)\le\delta_i$, and let $E\subseteq\R^n$ be a Borel
 set. Suppose that for every $i$ and every choice of real values of the
 coordinates other than the $i$th, the section of $E$ in the $i$th coordinate
 is a union of at most $C$ intervals (open, closed, half-open, degenerate, or
 unbounded). Then
 \[
  \Bigl|\Bigl(\bigotimes_i\mu_i\Bigr)(E)-\Bigl(\bigotimes_i\nu_i\Bigr)(E)\Bigr|
  \le2C\sum_{i=1}^n\delta_i .
 \]
\end{enumerate}
\end{lemma}

The section bound must hold for \emph{every} choice of the other
coordinates, not only for almost every choice under the proxy law, because
the replacement passes through mixed products of finite and continuous
marginals.

\begin{lemma}[A bounded rational Gaussian-like sampler]\label{lem:count:gauss-sampler}
For every integer $b\ge1$ there is a finite law $\mathcal G_b$ on $\Q$,
determined by $b$ through a fixed deterministic rule, with the following
properties.
\begin{enumerate}[label=(\alph*)]
\item Its support lies in $[-(b+20),b+20]$, and every atom is a rational
 number of bit length $O(b)$.
\item $d_{\rm K}(\mathcal G_b,N(0,1))\le2^{-b}$.
\item One sample is produced from independent fair bits in $\poly(b)$ bit
 operations in the worst case.
\end{enumerate}
\end{lemma}

The law $\mathcal G_{b,\sigma}$ is not a Gaussian law, and no theorem in this
paper takes real Gaussian coefficients as Turing input. The Gaussian law
$N(0,\sigma^2)$ serves only as a proxy in proofs: it makes orthogonal
projections of the noise independent and supplies density bounds. Every
theorem that uses a Gaussian-like law specifies the accuracy $b$ as a
function of the base data, computed before sampling. The bounds do not
apply to an arbitrary finite approximation of a Gaussian, and in particular
not to a coarse one, because Example~\ref{ex:count:atom} shows how atoms can
create unbounded counts.

\subsection{Global growth and active gradients}\label{sec:count:growth}

\begin{definition}[Point growth]\label{def:count:growth}
Let $X\subset\R^n$ be nonempty and compact, let $f:X\to\R$ be lower
semicontinuous, and let $F_c=f+c^{\mathsf T} x$. The \emph{point-growth modulus}
$g_*(c)\in[0,\infty]$ is the supremum of the numbers $g\ge0$ for which some
minimizer $x^*$ of $F_c$ on $X$ satisfies
\[
 F_c(x)-F_c(x^*)\ge g\norm{x-x^*}^2\qquad(x\in X).
\]
If $X$ is a singleton, $g_*=\infty$. If $F_c$ has two distinct minimizers,
then $g_*=0$. Write $w_i=\max_Xx_i-\min_Xx_i$ and $S=\sum_iw_i$.
\end{definition}

Point growth is measured from one minimizer. It differs from growth measured
by the distance to the set of minimizers, which can remain positive in the
presence of ties.

\begin{theorem}[Sharp growth tail]\label{thm:count:growth-tail}
Let $X\subset\R^n$ be nonempty and compact, let $f:X\to\R$ be lower
semicontinuous, and let $c$ have independent coordinates with densities
bounded by $\phi_1,\dots,\phi_n$. Then for every $\varepsilon>0$
\[
 \Prob\{g_*(c)<\varepsilon\}\le2\varepsilon\sum_{i=1}^n\phi_iw_i .
\]
In particular, if $c_i$ is uniform on $[-\sigma,\sigma]$ for every $i$, then
$\Prob\{g_*(c)<\varepsilon\}\le\varepsilon S/\sigma$. If $X$ is not a
singleton, then for every $\rho\in(0,1]$, with probability at least
$1-\rho$ the minimizer of $F_c$ is unique and
$g_*(c)\ge\rho/(2\sum_i\phi_iw_i)$. The constant $2$ cannot be decreased:
for $X=[-w/2,w/2]$, $f=0$, and $c$ uniform on $[-\sigma,\sigma]$,
$\Prob\{g_*(c)<\varepsilon\}=2\varepsilon\phi w$ with $\phi=1/(2\sigma)$
whenever $0<\varepsilon\le\sigma/w$.
\end{theorem}

The theorem assumes neither convexity, differentiability, nor definability of
$f$ or $X$. Its proof works in coefficient space. The function
$H(a)=\max_{x\in X}\{a^{\mathsf T} x-f(x)+\varepsilon\norm{x}^2\}$ is convex, and
wherever $H$ is differentiable at the proximal point of $-c$, the tilted
problem has growth at least $\varepsilon$. The proximal map is a
$1$-Lipschitz gradient map, so the area formula shows that its Jacobian is
singular almost everywhere on the bad event. This bounds the indicator of
the bad event by the trace of the Jacobian of the proximal residual, whose
integral along each coordinate is at most $2\varepsilon w_i$.

Two comparisons clarify the scope. For a finite binary feasible set and
independent continuous coefficients, Beier and V\"ocking bound the density of
the gap between the best and second-best values by the same expression
$2\sum_i\phi_i$ \cite{beier2006-typical-properties-of-winners-and}. That
statement compares values, not values relative to squared distances, and it
does not apply to continuous feasible sets. Qualitative genericity,
stability, and growth results are known for semialgebraic programs
\cite{lee2017-generic-properties-for-semialgebraic-programs}; they do not
give this quantitative tail on arbitrary compact sets with a lower
semicontinuous objective. Theorem~\ref{thm:count:growth-tail} quantifies the
probability of small growth.

The theorem concerns continuous laws. It does not by itself give an exact
algorithm in the bit model. For example, minimizing $x^2/2+cx$ over $[-1,1]$
with continuous $c$ has the interior minimizer $-c$, which is irrational with
probability one when it is interior. A finite rational law has atoms, and
ties can occur with positive probability. The next lemma gives finite-law
tails when the problem is semialgebraic.

For the polynomial statements, let $\mathcal D\subset\R^{n+m}$ be a nonempty
compact set described by a quantifier-free formula with at most $s$ atomic
polynomial sign conditions of degree at most $d$ and rational coefficients.
Let $F\in\Q[x,y]$ have degree at most $d$. The \emph{reduced objective}
$f(x)=\min\{F(x,y):(x,y)\in\mathcal D\}$ is lower semicontinuous on the
compact projection $X$ of $\mathcal D$. The noise acts on $x$ only:
$F_\gamma(x,y)=F(x,y)+\gamma^{\mathsf T} x$, and $g_*(\gamma)$ is the point-growth
modulus of $f+\gamma^{\mathsf T} x$ on $X$. When $m=0$, $f=F$. A bounded integer
coordinate $x_j\in\{\ell_j,\dots,u_j\}$ is described by the disjunction
$\bigvee_{t=\ell_j}^{u_j}(x_j=t)$, which adds $u_j-\ell_j+1$ atoms but no
variables.

\begin{lemma}[Finite-law growth and active-gradient tails]\label{lem:count:finite-tails}
There is an absolute constant $a_{\rm R}$ with the following properties. Put
\[
 H_{\rm s}=\bigl((2s+1)\max\{d,2\}\bigr)^{a_{\rm R}(n+m+1)^2},
 \qquad
 C_{\rm tail}=2H_{\rm s}^3+1 .
\]
\begin{enumerate}[label=(\alph*)]
\item \emph{(Sections.)} Fix $\varepsilon>0$, an index $i\le n$, and real
 values of $\gamma_l$ for $l\ne i$. The set of $t\in\R$ for which
 $g_*(\gamma)<\varepsilon$ when $\gamma_i=t$ is a union of at most
 $C_{\rm tail}$ intervals. The bound is uniform in $\varepsilon$ and in the
 fixed values.
\item \emph{(Uniform grids.)} If $\gamma_1,\dots,\gamma_n$ are independent
 with law $U_{\sigma,M}$, then for every $\varepsilon>0$
 \[
  \Prob\{g_*<\varepsilon\}\le\frac{S\varepsilon}{\sigma}+\frac{2nC_{\rm tail}}M,
  \qquad
  \Prob\{g_*=0\}\le\frac{2nC_{\rm tail}}M .
 \]
\item \emph{(Gaussian-like laws.)} If $\gamma_1,\dots,\gamma_n$ are
 independent with $d_{\rm K}(\gamma_i,N(0,\sigma^2))\le\delta_{\rm K}$, then
 for every $\varepsilon>0$
 \[
  \Prob\{g_*<\varepsilon\}\le\sqrt{2/\pi}\,\frac{S\varepsilon}{\sigma}
  +2nC_{\rm tail}\delta_{\rm K}.
 \]
\item \emph{(Active gradients on mixed boxes.)} Let $m=0$ and let
 $\mathcal D=X$ be a product of $n_c$ closed intervals of positive length
 and of bounded integer ranges with $R_Z$ integer points in total. Put
 $D=\max\{1,d-1\}$ and $K_{\rm act}=n_c3^{n_c}R_ZD^{n_c}$. For $\tau>0$, the
 probability that $g_*>0$ and the minimizer $x^*$ has a continuous
 coordinate $x_i^*$ at an endpoint of its interval with
 $|\partial_iF_\gamma(x^*)|\le\tau$ is at most
 $K_{\rm act}(\tau/\sigma+1/M)$ under the law of (b), and at most
 $K_{\rm act}(\sqrt{2/\pi}\,\tau/\sigma+2\delta_{\rm K})$ under the laws
 of (c).
\end{enumerate}
\end{lemma}

The constant $a_{\rm R}$ comes from the block-sensitive quantifier-elimination
theorem of Renegar \cite{renegar1992-on-the-computational-complexity-and}.
Its exponent depends on the product of the two block sizes, not on a tower
over the number of variables; a bound of cylindrical-decomposition type would
not give a sampling precision of polynomial length. For fixed $d$, if
$\log(s+1)$ is polynomial in $I$, then $\log C_{\rm tail}$ is polynomial in
$I$. This holds for mixed boxes with binary-encoded integer ranges, since
those contribute $\sum_j(u_j-\ell_j+1)$ atoms. The value of $a_{\rm R}$ is not
extracted here; the sampling rules below use one fixed admissible value. The
section bound is used only through the logarithm of $C_{\rm tail}$, which
determines a number of sampled bits. For quadratic objectives with $m=0$ an
elementary argument gives the explicit bound $8(\mathsf F+1)^2$ in place of
$C_{\rm tail}$, where $\mathsf F$ counts candidate faces
(Lemma~\ref{lem:qp:growth-sections}).

\begin{lemma}[Paying for a rare fallback]\label{lem:count:rare-fallback}
\begin{enumerate}[label=(\alph*)]
\item \emph{(Accounting.)} Let an algorithm act on one draw from a law fixed
 before sampling, and let its work on a draw be $W_0+\ind_{\mathcal E}W_1$.
 Suppose that $W_1\le B\,P(L_{\max})$ on every draw, where $B\ge1$ and the
 polynomial $P$ are fixed before sampling and $L_{\max}$ bounds the bit length
 of every sampled instance, and that $\Prob(\mathcal E)\le1/B$. Then
 $\E[W_0+\ind_{\mathcal E}W_1]\le\E W_0+P(L_{\max})$.
\item \emph{(Budget for polynomial mixed boxes.)} In the setting of
 Lemma~\ref{lem:count:finite-tails}(d), assume $S>0$, let $\rho\in(0,1]$, and
 put
 \[
  g_0=\frac{\rho\sigma}{S},\qquad
  \tau=\frac{\rho\sigma}{\max\{1,K_{\rm act}\}} .
 \]
 Let $\mathcal E_{\rm bad}$ be the event that $g_*<g_0$, or that $g_*\ge g_0$
 and some active continuous coordinate has $|\partial_iF_\gamma(x^*)|\le\tau$.
 If $\gamma_i\sim U_{\sigma,M}$ with
 $M\ge(2nC_{\rm tail}+K_{\rm act})/\rho$, or if
 $d_{\rm K}(\gamma_i,N(0,\sigma^2))\le\rho/(2nC_{\rm tail}+2K_{\rm act})$,
 then $\Prob(\mathcal E_{\rm bad})\le3\rho$.
\end{enumerate}
\end{lemma}

If $S=0$, the mixed box is a single point. It is evaluated directly before
any threshold is chosen, and neither the tails nor a fallback are needed;
part (b) divides by $S$ and is used only when $S>0$.

Part (a) cancels the combinatorial multiplier $B$ of a fallback, not its
whole bit cost: the remaining polynomial factor must depend only on the bit
length of the sampled instance, which is bounded before sampling. The
multiplier $B$ must cover every base-only cost of the fallback, including the
conversion of its output to the required format
(Theorem~\ref{thm:count:fallback}). This ordering, first $B$, then the
thresholds and the level cap, and only then the law, avoids a circular choice
in which a precision depends on the sampled coefficients that it is supposed
to control.

\subsection{An exact fallback for fixed-degree polynomial problems}\label{sec:count:fallback}

\begin{definition}[Explicit mixed semialgebraic instances]\label{def:count:semialg}
Fix a degree bound $d\ge1$. A \emph{base instance} consists of variables
$x\in\R^N$ with $N\ge1$, a set $\mathcal Z\subseteq\{1,\dots,N\}$ of integer variables
with binary-encoded bounds $\ell_j\le u_j$ ($j\in\mathcal Z$), a
quantifier-free formula $\Phi_0$ whose atoms are explicitly listed polynomial
sign conditions of degree at most $d$ with rational coefficients, and an
objective $F_0\in\Q[x]$ of degree at most $d$, given by its list of monomials
with rational coefficients. Its domain is
\[
 \mathcal D=\{x\in\R^N:\Phi_0(x)\text{ holds and }
 x_j\in\Z\cap[\ell_j,u_j]\text{ for }j\in\mathcal Z\},
\]
which is assumed nonempty and compact. Its bit length $I$ counts $\Phi_0$,
$F_0$, and the bounds. A \emph{tilted instance} adds rational coefficients
$\gamma_i$, each of bit length at most $b$, to any subset of the variables:
$F_\gamma(x)=F_0(x)+\sum_i\gamma_ix_i$. A real algebraic number $\theta$ is
given in \emph{root representation} by a nonzero $p\in\Z[t]$ and a rational
interval, possibly a point, that contains $\theta$ and no other real root of
$p$.
\end{definition}

The class contains mixed boxes ($\Phi_0$ a conjunction of bounds), bounded
polytopes and mixed polytopes, and polynomial graph constraints $y=g(x)$
with explicit polynomials $g$. It does not contain objectives given by
arithmetic circuits or by unexpanded compositions; the expanded monomial
list of $F_0$ is part of $I$. A problem with no remaining variables is
evaluated directly before this $N\ge1$ elimination construction is invoked.

\begin{theorem}[Exact fallback]\label{thm:count:fallback}
Fix $d\ge1$. There are a deterministic algorithm, a base-computable integer
$B=2^{\poly_d(I)}$, and an exponent $e_d$ depending only on $d$ such that the
following holds for every base instance of Definition~\ref{def:count:semialg}
and every tilted instance whose added coefficients have bit length at most
$b$.
\begin{enumerate}[label=(\roman*)]
\item Within $B\,(I+b+1)^{e_d}$ bit operations the algorithm returns an
 algebraic output
 (Definition~\ref{def:model:outputs}\ref{it:model:algebraic}) of the
 lexicographically least global minimizer $x^{\rm lex}$ and of the optimal
 value $f^*=\min_{\mathcal D}F_\gamma$: the integer coordinates of
 $x^{\rm lex}$; a squarefree polynomial $P\in\Z[t]$ with a rational interval
 isolating one real root $\theta$ of $P$; and rational polynomials $r_i$, for
 the continuous coordinates $i\notin\mathcal Z$, and $r_f$ with
 $x^{\rm lex}_i=r_i(\theta)$ and $f^*=r_f(\theta)$. As auxiliary data it also
 returns a root representation of every coordinate of $x^{\rm lex}$ and of
 $f^*$. The degree of $P$ and of every auxiliary defining polynomial is at
 most $B$, independently of $b$, and every coefficient and interval endpoint
 has bit length at most $B\,(I+b+1)^{e_d}$.
\item For every integer $q\ge0$ it returns, within $B\,(I+b+q+1)^{e_d}$ bit
 operations, a rational point within Euclidean distance $2^{-q}$ of
 $x^{\rm lex}$ and a rational interval of width at most $2^{-q}$ that
 contains $f^*$.
\item If $\mathcal D$ is a mixed box, it also returns within the bound of (ii)
 a rational point $y\in\mathcal D$ with $\norm{y-x^{\rm lex}}\le2^{-q}$ and a
 rational number $\ell\le f^*$ with $F_\gamma(y)-\ell\le2^{-q}$.
\end{enumerate}
Neither $B$ nor $e_d$ depends on $\gamma$, $b$, or $q$. The algorithm is
correct for every $\gamma$, including draws with ties, with a continuum of
minimizers, or with singular stationary sets.
\end{theorem}

The output is long on some instances: its length is at most
$B\,(I+b+1)^{e_d}$, exponential in the base length and polynomial in the
added coefficient length $b$. No bound independent of $b$ is possible, since
for the domain $\{1\}$ and $F_0=0$ the optimal value is the added coefficient
itself. Lemma~\ref{lem:count:rare-fallback}(a) needs exactly this separation
of a base-only factor from a polynomial factor in $I+b$. The proof
characterizes each coordinate of $x^{\rm lex}$, and the optimal value, as the
unique solution of a formula with two quantified blocks and one free
variable, eliminates the quantifiers by Renegar's theorem
\cite{renegar1992-on-the-computational-complexity-and}, and recovers each
singleton by univariate root isolation. The resulting scalar representations
all refer to the same canonical minimizer, but separately isolated
coordinates do not yet form an algebraic output. A final step converts them
into one shared root (Lemma~\ref{lem:count:shared-root}): in the tensor
product of the univariate quotient algebras, a linear form with integer
coefficients separates all tuples of roots; the characteristic polynomial $P$
of its multiplication matrix expresses every coordinate as a polynomial
function of its roots; and the scalar isolators select the root $\theta$
that corresponds to $x^{\rm lex}$ and $f^*$. The other roots of $P$
correspond to tuples of scalar roots that need not come from any optimizer.
Only the selected root carries meaning: the value identity
$F_\gamma(x^{\rm lex})=r_f(\theta)$ holds at $\theta$, but it need not hold as
an identity modulo $P$. The conversion costs a polynomial in the product of
the scalar degrees. Elimination bounds these degrees independently of $b$ (only
coefficient lengths grow with $b$), so the product is $2^{\poly_d(I)}$ and the
conversion fits in the factor $B$. That factor is fixed, with the conversion
included, before any threshold or law is chosen. No minimal polynomial,
unique minimizer, or finite stationary set is required. For quadratic
objectives on mixed polytopes, a much smaller rational fallback suffices
(Lemma~\ref{lem:qp:face}).

\subsection{How the tools are combined}\label{sec:count:template}

The closure-based exact results follow one template: those of
Section~\ref{sec:qp} and, with growth, label, or face certificates in place
of critical regions, those of Sections~\ref{sec:sp}--\ref{sec:int} other than
the two results described at the end of this subsection.
\begin{enumerate}[label=(\arabic*)]
\item \emph{Before sampling,} the algorithm computes from the base data a
 fallback multiplier $B$ that covers every base-only cost of the fallback,
 the bad-event thresholds, a level cap $J$, and a finite law (an accuracy $M$
 or $b$ and, for Gaussian-like laws, a support radius). All have polynomial
 bit length in $I$, except in the results with parameter-dependent sampling
 precision.
\item It draws the perturbation once.
\item It runs curvature-corrected refinement with exact closure steps through
 level $J$. Closure certifies the exact minimum of a cell from a structural
 formula valid on the whole cell. Correctness of every retained bound and of
 every closed cell holds on every draw (Theorem~\ref{thm:count:cells}).
\item If unresolved cells remain at level $J$, it runs an exact fallback on
 the same draw.
\end{enumerate}
The expected work of step (3) is bounded through
Theorem~\ref{thm:count:cells}(iv) by expected local counts of fixed grids,
computed under a continuous proxy and transferred by
Lemma~\ref{lem:count:transfer} when necessary. The probability that step
(4) runs is at most $1/B$, which pays for the fallback by
Lemma~\ref{lem:count:rare-fallback}(a). The structural routes differ in what
plays the role of the box $X$ and its function $f$, in the closure
certificate, and in the event whose small probability bounds step (4):
fixed hyperplane tubes and integer-label gaps for low-rank quadratic
problems (Section~\ref{sec:qp}), and growth and active-gradient events when
closure is certified by growth. Two results complete the sampled problem
differently and need no rare fallback. The lattice theorem
(Theorem~\ref{thm:int:lowrank}) stops at a predetermined level at which the
certified gap is below the spacing of the attainable objective values. The
strong-noise theorem (Theorem~\ref{thm:int:strong-field}) pins coordinates by
derivative signs and solves the remaining components exactly on every draw.
In each case the target is the exact optimum of the \emph{sampled} objective.

\paragraph{Laws that depend only on the input length.}
Each theorem computes its law from the individual base instance. In the
results with polynomial sampling precision one resolution per input length
suffices, because the base-only quantities that determine the law have
logarithms bounded by a fixed polynomial in the input length and satisfy the
schedule conditions below. To state this, fix a theorem together with its
fixed format constants (such as the degree and the dependency depth of graph
constraints) and its oracle implementation, let $\mathcal C$ be its class of
valid base instances, and write $I_\Pi$ for the base length of
$\Pi\in\mathcal C$. A \emph{grid schedule} assigns to every $\Pi$ a power of
two $M_\Pi$. It is \emph{monotone} if, for every power of two $M\ge M_\Pi$,
the conclusions of the theorem hold for $\Pi$ with $U_{\sigma,M}$ in place of
$U_{\sigma,M_\Pi}$, all other base-selected quantities of $\Pi$ (thresholds,
level cap, auxiliary box) unchanged, and the polynomial factor of the work
bound replaced by a fixed polynomial in $I_\Pi+\log_2M$. A \emph{Gaussian
support schedule} computes from $\Pi$, for each trial support radius
$2^t\sigma$ ($t=0,1,\dots$), an auxiliary box, a level cap, and a required
accuracy $b_\Pi(t)$, such that for every pair $(t,b)$ with $b\ge b_\Pi(t)$ and
$b+20\le2^t$ the conclusions of the theorem hold with $\mathcal G_{b,\sigma}$
and the box and cap of trial $t$, the polynomial factor being a fixed
polynomial in $I_\Pi+t+b$.

\begin{corollary}[A resolution depending only on the input length]\label{cor:count:universal-law}
Fix a theorem as above and an integer $I\ge1$.
\begin{enumerate}[label=(\alph*)]
\item\label{it:count:ul-grid} Suppose its grid schedule is monotone and
 $\log_2M_\Pi\le P_0(I_\Pi)$ for all $\Pi\in\mathcal C$, with an effective
 nondecreasing polynomial $P_0\ge1$. Then for every $\Pi\in\mathcal C$ with
 $I_\Pi\le I$ the theorem holds with the law $U_{\sigma,M(I)}$,
 $M(I)=2^{\lceil P_0(I)\rceil}$. The perturbation model, the scale $\sigma$
 supplied with $\Pi$, the other base-selected quantities of $\Pi$,
 correctness on every draw, the output contract, and the structural and
 numerical factors of the expected work bound are unchanged; the polynomial
 factor becomes a polynomial in $I$ whose exponent depends only on the
 exponents of $P_0$ and of the original factor.
\item\label{it:count:ul-gauss} Suppose it has a Gaussian support schedule
 with $b_\Pi(t)\le A(I_\Pi)+D(I_\Pi)\,t$ for all $\Pi\in\mathcal C$ and
 $t\ge0$, where $A,D\ge1$ are effective nondecreasing integer polynomials, and
 put
 \[
  H(I)=A(I)+D(I)+21,\qquad t(I)=\lceil4\log_2H(I)\rceil,\qquad
  b(I)=A(I)+D(I)\,t(I).
 \]
 Then for every $\Pi\in\mathcal C$ with $I_\Pi\le I$ the theorem holds with
 $\mathcal G_{b(I),\sigma}$ and the box and cap of trial $t(I)$, with the
 same conclusions as in (a). The support radius $(b(I)+20)\sigma$ is at most
 $2^{t(I)}\sigma\le2H(I)^4\sigma$.
\item\label{it:count:ul-param} Suppose that the hypotheses of (a) hold except
 that $\log_2M_\Pi\le F(K)\,P_0(I_\Pi)$ is known only for the instances whose
 structural parameter is at most a supplied bound $K$, where $F$ is an
 nondecreasing positive integer-valued function, computable in
 $(1+\operatorname{size}(K)+F(K))^{c_F}$ bit operations for a fixed
 exponent $c_F$. The supplied bound $K$ is included in the base data.
 Then (a) holds for these instances with
 $M(I,K)=2^{\lceil F(K)P_0(I)\rceil}$, except that the bounds on the sampled
 coefficient lengths, on the output length, and the polynomial factor of the
 work bound are multiplied by $(1+F(K))^{c}$, with $c$ depending only on the
 exponents of $P_0$, of the original bounds, and $c_F$.
\end{enumerate}
\end{corollary}

The hypotheses hold route by route as follows; the proofs are in
Appendix~\ref{app:count:universal}. The closure-based grid theorems of
Sections~\ref{sec:qp}--\ref{sec:int} with polynomial sampling precision,
and Theorem~\ref{thm:qp:two}(a), choose $M_\Pi$ as the least power of two
above finitely many
base-only quantities, such as $2^J$ and inverse failure budgets, whose
logarithms are bounded by a fixed polynomial in $I_\Pi$. It uses $M$ only
through the interval masses \eqref{eq:model:interval}, the replacement errors
of Lemma~\ref{lem:count:transfer}, and these lower bounds, so its schedule is
monotone and part (a) applies; for the quadratic theorems the verification is
in Appendices~\ref{app:qp:main} and~\ref{app:qp:sep}. The lattice theorem
(Theorem~\ref{thm:int:lowrank}) is the one exception to the literal
definition, because its level cap is $J=\log_2M$; part (a) applies to it once
the cap is reset to $\log_2M(I)$. Part (b) applies to the Gaussian-like
Theorems~\ref{thm:qp:two}(b), \ref{thm:qp:gauss}, and~\ref{thm:qp:sep-gauss}.
Part (c) applies to the core-only flow and totally unimodular theorems with
arbitrary core faces (Theorems~\ref{thm:int:flow-boundary}
and~\ref{thm:int:tu}), whose sampling precision depends on the core size $k$,
with a supplied $K\ge k$. Taking $K=I$ gives a law that depends only on $I$,
but its sampling cost need not stay within the bound for the actual core
size; nothing here gives these two theorems a law of polynomial sampling
precision. The strong-noise theorem holds for every $M$, but its bound
depends on $M$ through probabilities $q_i$ that are not monotone in $M$: the
grids of sizes $2$ and $4$ give the interval $[-\sigma/2,\sigma/2]$ the masses
$0$ and $1/2$. With a common $M$, its $q_i$ and $\beta$ must be computed for
that $M$; its sufficient regime \eqref{eq:int:sf-regime}, a lower bound on
$\sigma$ together with a lower bound on $M$, is preserved when $M$ grows.

Three qualifications apply. First, the hypothesis
$\log_2M_\Pi\le P_0(I_\Pi)$ is a property of the domains of a route, not of
fixed degree alone. The fallback factor and the tail constants of
Lemma~\ref{lem:count:finite-tails} have logarithms polynomial in $I$ for
every explicit semialgebraic instance, since they depend on the number $s$ of
atoms only through $\log(s+1)$, even when native integer ranges contribute
exponentially many atoms. The level caps also involve widths and curvature
bounds, and on general compact semialgebraic domains these need not have
polynomial bit length: the domain $1\le x_0\le2$, $x_i=x_{i-1}^2$
($1\le i\le N$) has a description of length $O(N\log N)$ and last width
$2^{2^N}-1$. The theorems above use boxes, rational polytopes, explicit graph
constraints of fixed dependency depth, and implicit graphs with supplied
brackets, on which these quantities have polynomial encoding length. Second,
the oracle contract of a theorem must hold for every coefficient in the
support of the new law, as the contracts in this paper do, and oracle costs
remain explicit factors. Third, the common law standardizes only the scalar
law: it is rescaled by the supplied $\sigma$ and applied to the coordinates
that the theorem perturbs, so it is not one probability measure for all
instances, and grid and Gaussian-like laws are not interchangeable. For
Gaussian-like laws the support radius $(b(I)+20)\sigma$ is the one that
enters the original-objective bound of Proposition~\ref{prop:model:regret}.
Arbitrary-precision evaluation of an output uses the same draw and does not
change the law; only Corollary~\ref{cor:count:approx}, whose cap depends on
the requested accuracy, needs a resolution depending on that accuracy. The
corollary is an elementary consequence of the effective bounds used in the
proofs. A finite collection of routes can share an enlarged envelope within
each law family when its format constants and oracle cost models are fixed.
